% Part: intuitionistic-logic
% Chapter: tableaux
% Section: introduction

\documentclass[../../../include/open-logic-section]{subfiles}

\begin{document}

\olfileid{int}{tab}{int}

\olsection{مقدمه}

!!^{tableau}ها درخت‌هایی معین (با انشعاب رو به پایین) از !!{signed
  formula}ها هستند؛ یعنی جفت‌هایی متشکل از نشانهٔ ارزش صدق
($\True$ یا~$\False$) و !!a{formula}:
\[
\sFmla{\True}{!A} \text{ یا } \sFmla{\False}{!A}.
\]
!!^a{tableau} با شماری \emph{فرض} آغاز می‌شود. هر !!{signed formula}
بعدی با اعمال یکی از قواعد استنتاج تولید می‌شود. برخی قواعد استنتاج یک
یا چند !!{signed formula} به انتهای درخت می‌افزایند؛ برخی دیگر دو
انتهای تازه می‌افزایند و دو شاخه پدید می‌آورند. قواعد به !!{signed
formula}هایی می‌انجامند که !!{formula} آن‌ها از !!{signed formula}ای
که قاعده بر آن اعمال شده ساده‌تر است. هرگاه شاخه‌ای هم
$\sFmla{\True}{!A}$ و هم~$\sFmla{\False}{!A}$ را دربرداشته باشد،
می‌گوییم شاخه \emph{بسته} است. اگر همهٔ شاخه‌های !!a{tableau} بسته
باشند، کل !!{tableau} بسته است. یک !!{tableau} بسته، !!a{derivation}ای
است که نشان می‌دهد مجموعهٔ !!{signed formula}هایی که !!{tableau} با
آن‌ها آغاز شده ارضاناپذیر است. می‌توان از این امر برای تعریف رابطهٔ
$\Proves$ استفاده کرد: $\Gamma \Proves !A$ اگر و تنها اگر مجموعهٔ
متناهی~$\Gamma_0 = \{!B_1, \dots, !B_n\} \subseteq \Gamma$ای وجود
داشته باشد چنان‌که یک !!{tableau} بسته برای فرض‌های زیر وجود داشته باشد:
\[
\{\sFmla{\False}{!A}, \sFmla{\True}{!B_1}, \dots, \sFmla{\True}{!B_n}\}.
\]

برای منطق شهودگرایانه باید هم مفهوم !!{signed formula} را گسترش دهیم
و هم قواعد رابط‌ها را تعدیل کنیم. افزون بر نشانه ($\True$ یا~$\False$)،
!!{formula}ها در !!{tableau}های شهودگرایانه \emph{پیشوندهای}
$\sigma$ نیز دارند. پیشوندها دنباله‌های ناتهی از اعداد صحیح مثبت‌اند؛
یعنی~$\sigma \in (\PosInt)^* \setminus \{\emptyseq\}$. این پیشوندها
را بدون~$\tuple{\ }$ پیرامونی می‌نویسیم و !!{element}های منفرد را به‌جای ویرگول
با نقطه جدا می‌کنیم. اگر~$\sigma$ پیشوندی باشد، آنگاه
$\sigma.n$ همان~$\sigma \concat \tuple{n}$ است؛ برای نمونه، اگر
$\sigma = 1.2.1$، آنگاه~$\sigma.3$ برابر~$1.2.1.3$ است. پس، مثلاً،
\[
\sFmla{\True}{!A \lif (!B \lif !C)}[1.2]
\]
یک \emph{!!{signed formula} پیشونددار} (یا به‌اختصار یک
\emph{!!{formula} پیشونددار}) است.

به‌طور شهودی، پیشوند نام جهانی در مدلی است که ممکن است !!{formula}های
یک شاخه از !!a{tableau} را ارضا کند؛ و اگر~$\sigma$ نام جهانی باشد،
$\sigma.n$ نام جهانی است که از (جهان موسوم به)~$\sigma$ دسترس‌پذیر است.

در مدل‌های شهودگرایانه رابطهٔ دسترسی‌پذیری بازتابی و تعدی است. به زبان
پیشوندها، یعنی~$\sigma$ از خودش دسترس‌پذیر است و هر پیشوندی که
$\sigma$ را امتداد دهد، یعنی هر پیشوندِ صورت
$\sigma.n_1.\cdots.n_k$، نیز از آن دسترس‌پذیر است. نماد~$\sigma.*$ را
برای نشان‌دادن خود~$\sigma$ و هر امتداد آن معرفی می‌کنیم. به بیان دیگر،
پیشوندهای~$\sigma.*$ دقیقاً همهٔ پیشوندهای دسترس‌پذیر از~$\sigma$ هستند.

\end{document}
